\begin{helper}[A second tree vertex determines last-return data]
Fix a defective color $\delta$.  If $w\in V(T)$, $w\ne v^0_\delta$, and
the vertex-node $w$ belongs to the full
$(\gamma(\delta),\delta)$ component through $v^0_\delta$, then there
exists a \noderef{ScheideDefectiveReturnData} for $\delta$.
\end{helper}
\begin{proof}
Write $X=\noderef{TreeSequenceVertexSet}(T)$,
$i=\gammaIndex(\delta)$, $\gamma=\gamma_i$, and let $P'$ be the
replacement path supplied by \noderef{ScheideDefectiveSetup}.  The
predicate \noderef{ScheideDefectiveChainPath} recorded there says that
$P'$ is simple, starts at the tree vertex $v^0_\delta$, ends at a vertex
outside $X$, and has support exactly the full
$(\gamma,\delta)$ component through $v^0_\delta$.  The component
hypothesis therefore puts the vertex-node $w$ on $P'$.  Its occurrence
is not the initial occurrence, because $w\ne v^0_\delta$ and $P'$ is
simple.

Consider the finite set of positions of $P'$ occupied by a vertex-node
whose vertex lies in $X$, and choose its largest position $r$.  This set
contains the occurrence of $w$, so $r>0$.  Write $v_1$ for the vertex at
position $r$.  Simplicity then gives $v_1\ne v^0_\delta$.  Since the
terminal vertex of $P'$ lies outside $X$, one has $r<|P'|$.  By the
bipartite definition \noderef{TwoColorIncidenceGraph}, the node at
$r+1$ is an edge-node, say $f_1$.  It cannot be the terminal node, which
is a vertex-node, so the node at $r+2$ exists and is a vertex-node, say
$u_1$.  Maximality of $r$ proves both $u_1\notin X$ and, more generally,
that every vertex-node strictly after $r$ has its vertex outside $X$.
The two successive adjacencies give
$\noderef{Incident}(G,v_1,f_1)$ and
$\noderef{Incident}(G,u_1,f_1)$.  These are distinct vertices by
simplicity of $P'$.  Hence the two-element endpoint object of the
loopless edge $f_1$ is exactly
$\operatorname{ends}(f_1)=\{v_1,u_1\}$, and
\noderef{BoundaryEdges} gives $f_1\in\partial_G(X)$.

We next use the full path, not its suffix, to prove that $v_1$ is an
exit.  Interchanging the two displayed disjuncts in
\noderef{TwoColorIncidenceGraph} gives the identity graph homomorphism
from the $(\gamma,\delta)$ incidence graph to the
$(\delta,\gamma)$ incidence graph.  Mapping $P'$ through this
homomorphism changes neither its nodes, its support, its simplicity, nor
the chosen index $r$.  The reverse identity homomorphism shows that the
two connected components have equal support.  Thus the mapped full path
still satisfies \noderef{ScheideDefectiveChainPath}, now in the formal
color order $(\delta,\gamma)$.  Its initial vertex is
$v^0_\delta\in X$, the setup states that $\gamma$ is missing there,
$\delta$ is defective, and maximality of $r$ says that $v_1$ is its last
tree vertex.  These are precisely the witnesses required by
\noderef{ScheideExitVertex}.  Consequently
$v_1\in\noderef{ScheideExitVertices}(T,\varphi)$, and
\noderef{ScheideStructuralExitsInBase} places $v_1$ in the base image
$Y$.

The edge-node adjacency says, again by
\noderef{TwoColorIncidenceGraph}, that $f_1$ has color $\gamma$ or
$\delta$.  The inherited data in
\noderef{ScheideStructuralConfiguration}, invoked through
\noderef{ScheideReducedConfiguration}, include a globally maximum
triple.  Hence \noderef{ScheideGlobalMaximumOrder} says that $X$ is
\noderef{ColoredVertexSetClosed} and
\noderef{ColoredVertexSetElementary}.  Since $f_1$ is a boundary edge
and $\gamma$ is missing at $v^0_\delta\in X$, closedness rules out the
first color alternative.  Therefore $\varphi(f_1)=\delta$.

The exact-one exit clause of \noderef{ScheideReducedConfiguration}
gives a unique block index $j$ such that
$\gamma_j\in\noderef{MissingColors}(G,\varphi,v_1)$.  If
$\gamma_j=\gamma_i$, that color would be missing at both distinct tree
vertices $v_1$ and $v^0_\delta$, contradicting the just-established
elementariness of $X$.  Thus $\gamma_j\ne\gamma_i$; injectivity of the
block-color map in \noderef{ScheideStructuralConfiguration} also gives
$j\ne i$.

It remains to check the other color inequalities required by
\noderef{ScheideDefectiveReturnData}.  The setup states directly that
$\alpha _0$ differs from every block color, hence from both
$\gamma_i$ and $\gamma_j$.  The base-position equations in
\noderef{ScheideStructuralConfiguration} explicitly put $y_0$ in $X$;
only now do we use that $\alpha _0$ is missing at $y_0$.  Membership of
$\delta$ in \noderef{DefectiveColors}, together with the definition
\noderef{BoundaryColorEdges}, supplies at least one boundary identity
colored $\delta$.  By \noderef{ColoredVertexSetClosed}, this same color
is missing at no vertex of $X$.  Applying this fact at
$y_0,v^0_\delta,v_1$ shows respectively that
$\delta\ne\alpha _0$, $\delta\ne\gamma_i$, and
$\delta\ne\gamma_j$.  The positions $r,r+1,r+2$, the universal
later-vertex assertion, the endpoint and color identities, the exit and
base memberships, the missing-color witness at $v_1$, and these six
pairwise inequalities fill every field of the required return datum.
\end{proof}
